\documentclass[main.tex]{subfiles}
\begin{document}
\section{Expliciter les savoirs et les procédures}
%%%%%%%%%%%%%%%%%%%%%%%%%%%%%%%%%%%%%%%
% EXERCICE
\begin{exercise}[points=1]
    Invente une série statistique de 5 valeurs dont la médiane est 4.
\end{exercise}
% SOLUTION
\begin{solution}
    \numlist{1;2;4;5;6}
\end{solution}
%%%%%%%%%%%%%%%%%%%%%%%%%%%%%%%%%%%%%%%
% EXERCICE
\begin{exercise}[points=3]
    Pour chaque variable ci-dessous, détermine si elle est qualitative ou quantitative.
    \begin{center}
        \begin{tblr}{column{2,3}={c}}
                                     & Quantitative    & Qualitative     \\
            Les couleurs             & \faIcon{square} & \faIcon{square} \\
            La nationalité           & \faIcon{square} & \faIcon{square} \\
            Le nombre d'enfants      & \faIcon{square} & \faIcon{square} \\
            La pointure de chaussure & \faIcon{square} & \faIcon{square} \\
            La taille                & \faIcon{square} & \faIcon{square} \\
            Le sexe d'un individu    & \faIcon{square} & \faIcon{square}
        \end{tblr}
    \end{center}
\end{exercise}
% SOLUTION
\begin{solution}
    \begin{center}
        \begin{tblr}{column{2,3}={c}}
                                     & Quantitative    & Qualitative     \\
            Les couleurs             & \faIcon{square} & \faCheckSquare \\
            La nationalité           & \faIcon{square} & \faCheckSquare \\
            Le nombre d'enfants      & \faCheckSquare & \faIcon{square} \\
            La pointure de chaussure & \faCheckSquare & \faIcon{square} \\
            La taille                & \faCheckSquare & \faIcon{square} \\
            Le sexe d'un individu    & \faCheckSquare & \faIcon{square}
        \end{tblr}
    \end{center}
\end{solution}
%%%%%%%%%%%%%%%%%%%%%%%%%%%%%%%%%%%%%%%
% EXERCICE
\newpage
\begin{exercise}[points=4]
    Pour chaque situation, coche uniquement la réponse correcte.
    \begin{enumerate}
        \item Une municipalité veut connaître l'opinion de ses administrés sur le programme de réfection de la voirie et de l'éclairage public de la ville. La responsable de la communication a sélectionné au hasard 75 noms dans l'annuaire téléphonique de la ville pour sonder ces personnes par téléphone.
              \begin{QCM}
                  \choice La population est l'ensemble des administrés figurant dans l'annuaire ; l'échantillon est constitué des 75 personnes sélectionnées dans cet annuaire.

                  \choice La population est l'ensemble des administrés de la ville ; l'échantillon est constitué des administrés inscrits sur la liste électorale.

                  \choice La population est l'ensemble des administrés inscrits sur la liste électorale ; l'échantillon est constitué des personnes inscrites sur l'annuaire téléphonique de la ville.
              \end{QCM}
        \item Un contrôleur de qualité dans une usine automobile veut vérifier l'épaisseur de la peinture sur une voiture. Il choisit au hasard 30 zones sur la voiture et mesure l'épaisseur de la peinture dans chacune de ces zones.
              \begin{QCM}
                  \choice La population est l'ensemble des zones de la voiture ; l'échantillon est constitué des 30 zones sélectionnées.

                  \choice La population est l'ensemble des voitures produites ; l'échantillon est l'une de ces voitures.

                  \choice La population est l'ensemble des voitures produites ; l'échantillon est constitué des 30 zones sélectionnées.
              \end{QCM}
        \item Pour savoir si les élèves de Terminale du lycée étaient satisfaits de la cantine, le conseiller pédagogique d'éducation a interrogé 100 élèves de Terminale choisis au hasard.
              \begin{QCM}
                  \choice La population est l'ensemble des élèves de Terminale de France ; l'échantillon est l'ensemble des élèves de Terminale de ce lycée.

                  \choice La population est l'ensemble des élèves de ce lycée ; l'échantillon est l'ensemble des élèves de Terminale de ce lycée.

                  \choice La population est l'ensemble des élèves de Terminale de ce lycée ; l'échantillon est constitué de 100 élèves de Terminale de ce lycée.
              \end{QCM}
        \item Le propriétaire des stands de boisson d'un stade de football voudrait vendre une nouvelle boisson. Pour savoir quelle future boisson aura le plus de succès, il choisit au hasard 80 numéros de siège et interroge leurs occupants sur leur boisson préférée.
              \begin{QCM}
                  \choice La population est l'ensemble des visiteurs du stade qui achètent des boissons ; l'échantillon est constitué des 80 sièges sélectionnés.

                  \choice La population est l'ensemble des occupants des 80 sièges sélectionnés ; l'échantillon est constitué des différentes futures boissons.

                  \choice La population est l'ensemble des occupants des sièges du stade ; l'échantillon est constitué des occupants des 80 sièges sélectionnés.
              \end{QCM}
    \end{enumerate}
\end{exercise}
% SOLUTION
\begin{solution}
    \begin{enumerate}
        \item Une municipalité veut connaître l'opinion de ses administrés sur le programme de réfection de la voirie et de l'éclairage public de la ville. La responsable de la communication a sélectionné au hasard 75 noms dans l'annuaire téléphonique de la ville pour sonder ces personnes par téléphone.
              \begin{QCM}
                  \choice[\faIcon{check-square}] La population est l'ensemble des administrés figurant dans l'annuaire ; l'échantillon est constitué des 75 personnes sélectionnées dans cet annuaire.

                  \choice La population est l'ensemble des administrés de la ville ; l'échantillon est constitué des administrés inscrits sur la liste électorale.

                  \choice La population est l'ensemble des administrés inscrits sur la liste électorale ; l'échantillon est constitué des personnes inscrites sur l'annuaire téléphonique de la ville.
              \end{QCM}
        \item Un contrôleur de qualité dans une usine automobile veut vérifier l'épaisseur de la peinture sur une voiture. Il choisit au hasard 30 zones sur la voiture et mesure l'épaisseur de la peinture dans chacune de ces zones.
              \begin{QCM}
                  \choice[\faIcon{check-square}] La population est l'ensemble des zones de la voiture ; l'échantillon est constitué des 30 zones sélectionnées.

                  \choice La population est l'ensemble des voitures produites ; l'échantillon est l'une de ces voitures.

                  \choice La population est l'ensemble des voitures produites ; l'échantillon est constitué des 30 zones sélectionnées.
              \end{QCM}
        \item Pour savoir si les élèves de Terminale du lycée étaient satisfaits de la cantine, le Conseiller pédagogique d'Éducation a interrogé 100 élèves de Terminale choisis au hasard.
              \begin{QCM}
                  \choice La population est l'ensemble des élèves de Terminale de France ; l'échantillon est l'ensemble des élèves de Terminale de ce lycée.

                  \choice La population est l'ensemble des élèves de ce lycée ; l'échantillon est l'ensemble des élèves de Terminale de ce lycée.

                  \choice[\faIcon{check-square}]La population est l'ensemble des élèves de Terminale de ce lycée ; l'échantillon est constitué de 100 élèves de Terminale de ce lycée.
              \end{QCM}
        \item Le propriétaire des stands de boisson d'un stade de football voudrait vendre une nouvelle boisson. Pour savoir quelle future boisson aura le plus de succès, il choisit au hasard 80 numéros de siège et interroge leurs occupants sur leur boisson préférée.
              \begin{QCM}
                  \choice La population est l'ensemble des visiteurs du stade qui achètent des boissons ; l'échantillon est constitué des 80 sièges sélectionnés.

                  \choice La population est l'ensemble des occupants des 80 sièges sélectionnés ; l'échantillon est constitué des différentes futures boissons.

                  \choice[\faIcon{check-square}] La population est l'ensemble des occupants des sièges du stade ; l'échantillon est constitué des occupants des 80 sièges sélectionnés.
              \end{QCM}
    \end{enumerate}
\end{solution}
\end{document}